\documentclass[11pt,a4paper]{article}
\usepackage[utf8]{inputenc}
\usepackage[T1]{fontenc}
\usepackage[spanish,es-minimal,es-noshorthands]{babel}
\AtBeginDocument{\renewcommand{\sin}{\operatorname{sen}}\renewcommand{\arcsin}{\operatorname{arcsen}}}
\usepackage{lmodern}
\usepackage[sfdefault,lining]{FiraSans}
\usepackage[a4paper,top=2.3cm,bottom=2.2cm,left=1.8cm,right=1.8cm,headheight=1.3cm,headsep=0.4cm,footskip=1cm]{geometry}
\usepackage{amsmath,amssymb}
\usepackage{icomma}
\usepackage{graphicx}
\usepackage[table]{xcolor}
\usepackage{tikz}
\usepackage{fontawesome5}
\usepackage{ccicons}
\usepackage{fancyhdr}
\usepackage{lastpage}
\usepackage{tasks}
\usepackage{needspace}
\usepackage[most]{tcolorbox}
\usepackage{array}
\usepackage[hidelinks]{hyperref}
\AtBeginDocument{\hypersetup{pdftitle={Ficha de ejercicios de la unidad 6: Proporcionalidad y porcentajes. 2º ESO},pdfsubject={Matemáticas 2º ESO},pdfkeywords={razones, proporciones, proporcionalidad directa, proporcionalidad inversa, regla de tres, proporcionalidad compuesta, repartos proporcionales, porcentajes, aumentos y disminuciones porcentuales}}}

\definecolor{navy}{HTML}{1B2748}
\definecolor{accent}{HTML}{E8603A}
\definecolor{teal}{HTML}{1F9E90}
\definecolor{soft}{HTML}{EEF1F8}
\definecolor{muted}{HTML}{5D6685}

\setlength{\parindent}{0pt}
\setlength{\parskip}{3pt}
\settasks{label=\alph*), label-format=\bfseries\color{navy}, label-width=1.3em, item-indent=1.8em, column-sep=0.6em, after-item-skip=5pt, before-skip=3pt, after-skip=4pt}

% ---------- capçalera i peu ----------
\newcommand{\llicencia}{\href{https://creativecommons.org/licenses/by-nc-sa/4.0/deed.es}{\ccbyncsa}\ \ Licencia Creative Commons Atribución-NoComercial-CompartirIgual 4.0 Internacional (CC BY-NC-SA 4.0)}
\pagestyle{fancy}
\fancyhf{}
\fancyhead[L]{\raisebox{-0.35\height}{\includegraphics[height=0.95cm]{logo-gva.png}}}
\fancyhead[C]{\small\color{muted}Ficha de ejercicios · Unidad 6: Proporcionalidad y porcentajes · 2\textsuperscript{o} ESO}
\fancyhead[R]{\raisebox{-0.35\height}{\includegraphics[height=0.95cm]{logo-vermellar.png}}}
\fancyfoot[L]{\scriptsize\color{muted}\llicencia}
\fancyfoot[R]{\small\color{muted}\thepage}
\renewcommand{\headrulewidth}{0.4pt}
\renewcommand{\headrule}{\color{muted!40}\hrule height \headrulewidth}
\fancypagestyle{primera}{\fancyhf{}\renewcommand{\headrulewidth}{0pt}
  \fancyfoot[L]{\scriptsize\color{muted}\llicencia}
  \fancyfoot[R]{\small\color{muted}\thepage}}

% ---------- nivells i exercicis ----------
\newcommand{\estrelles}[1]{{\color{accent}\ifcase#1\or\faStar\faStar[regular]\faStar[regular]\or\faStar\faStar\faStar[regular]\or\faStar\faStar\faStar\fi}}
\newcounter{exer}
\newcommand{\ex}[1]{\par\needspace{4\baselineskip}\medskip\stepcounter{exer}\noindent{\color{navy}\bfseries\theexer.}\ \ {\small\estrelles{#1}}\ \ }

\newcommand{\seccio}[2]{\par\needspace{10\baselineskip}\bigskip
  \begin{tcolorbox}[colback=navy,colframe=navy,coltext=white,arc=2mm,boxrule=0pt,left=3mm,right=3mm,top=1.2mm,bottom=1.2mm]
  {\large\bfseries #1\quad #2}
  \end{tcolorbox}\vspace{-2pt}}

\newcolumntype{C}[1]{>{\centering\arraybackslash}p{#1}}
\usepackage{pifont}
% quadre de recordatori (només a la fitxa d'exercicis)
\newcommand{\recorda}[1]{\par\smallskip\begin{tcolorbox}[colback=accent!7,colframe=accent!55,arc=2mm,boxrule=0.5pt,left=2.5mm,right=2.5mm,top=1mm,bottom=1mm]\small{\color{accent}\faLightbulb[regular]}\ \textbf{Recuerda:}\ #1\end{tcolorbox}}
\newcommand{\res}[1]{#1}\newcommand{\rest}[1]{#1}

\def\fmneed{14}
% =====================================================================
% Capçalera FentMates de les fitxes: un color per a cada curs.
% Els colors són els de la pàgina inicial (LEVELS a assets/js/materials.js);
% el logo ve de _eines/logo/pdf/ (python3 _eines/logo/genera.py).
%
% Ús, al final del preàmbul de la unitat:
%   % =====================================================================
% Capçalera FentMates de les fitxes: un color per a cada curs.
% Els colors són els de la pàgina inicial (LEVELS a assets/js/materials.js);
% el logo ve de _eines/logo/pdf/ (python3 _eines/logo/genera.py).
%
% Ús, al final del preàmbul de la unitat:
%   % =====================================================================
% Capçalera FentMates de les fitxes: un color per a cada curs.
% Els colors són els de la pàgina inicial (LEVELS a assets/js/materials.js);
% el logo ve de _eines/logo/pdf/ (python3 _eines/logo/genera.py).
%
% Ús, al final del preàmbul de la unitat:
%   \input{../comu/fentmates.tex}
%   \fmcurs{1eso}
% i, al començament del document:
%   \capcalera{Matemàtiques}{1r ESO}{Fitxa d'exercicis · Unitat 2}{Divisibilitat}
% Opcional abans de \input: \def\fmneed{14} (espai mínim abans de cada secció).
% =====================================================================
\makeatletter
\def\fm@logodir{../../logo/pdf}
\providecommand{\fmneed}{10}
\def\fm@cursnom{}

% clau del curs, color i color del text damunt del color
\newcommand{\fm@curs}[3]{\@namedef{fm@color@#1}{#2}\@namedef{fm@text@#1}{#3}}
\fm@curs{1eso}{E6B84A}{navy}
\fm@curs{2eso}{3B82F6}{white}
\fm@curs{3eso}{FF7A50}{white}
\fm@curs{4eso-a}{A78BFA}{white}
\fm@curs{4eso-b}{6D28D9}{white}
\fm@curs{1batx-mi}{9B1C1C}{white}
\fm@curs{1batx-ccss}{EF6B6B}{white}
\fm@curs{2batx-mii}{1E6B47}{white}
\fm@curs{2batx-ccss}{2BB3A3}{white}

\newcommand{\fmcurs}[1]{%
  \@ifundefined{fm@color@#1}{\PackageError{fentmates}{Curs desconegut: #1}{}}{}%
  \edef\fm@tmp{\noexpand\definecolor{curs}{HTML}{\@nameuse{fm@color@#1}}%
               \noexpand\colorlet{curstext}{\@nameuse{fm@text@#1}}}\fm@tmp
  \colorlet{cursfosc}{curs!72!black}%
  \def\fm@logo{\fm@logodir/vertical-#1.pdf}%
  \def\fm@marca{\fm@logodir/marca-#1.pdf}}

% ---------- capçalera de la primera pàgina ----------
% #1 matèria (petit), #2 curs, #3 tipus de document i unitat, #4 títol
\newcommand{\capcalera}[4]{%
  \gdef\fm@cursnom{#2}%
  \par\noindent
  \begin{tcolorbox}[colback=curs,colframe=curs,boxrule=0pt,arc=2.2mm,width=3.5cm,height=2.15cm,
      valign=center,left=3mm,right=2mm,top=1mm,bottom=1mm,nobeforeafter,box align=center]
    \color{curstext}{\footnotesize #1}\par\vspace{1pt}{\fontsize{19}{21}\selectfont\bfseries #2\par}%
  \end{tcolorbox}\hspace{5mm}%
  \begin{minipage}[c]{\dimexpr\linewidth-3.5cm-5mm-2.9cm\relax}
    {\color{cursfosc}\small\bfseries\MakeUppercase{#3}}\par\vspace{2pt}
    {\color{navy}\fontsize{20}{23}\selectfont\bfseries #4\par}
  \end{minipage}\hfill
  \begin{minipage}[c]{2.6cm}\raggedleft\includegraphics[width=2.6cm]{\fm@logo}\end{minipage}\par
  \vspace{2pt}}

% ---------- barres de secció del color del curs ----------
\renewcommand{\seccio}[2]{\par\needspace{\fmneed\baselineskip}\bigskip
  \begin{tcolorbox}[colback=curs,colframe=curs,coltext=curstext,arc=2mm,boxrule=0pt,left=3mm,right=3mm,top=1.2mm,bottom=1.2mm]
  {\large\bfseries #1\quad #2}
  \end{tcolorbox}\vspace{-2pt}}

% ---------- estrelles de dificultat del color del curs (a Batxillerat no se'n posen) ----------
\@ifundefined{estrelles}{}{%
  \renewcommand{\estrelles}[1]{{\color{curs}\ifcase#1\or\faStar\faStar[regular]\faStar[regular]\or\faStar\faStar\faStar[regular]\or\faStar\faStar\faStar\fi}}}

% ---------- capçalera i peu de les pàgines ----------
% primera pàgina: logos institucionals al peu
\fancypagestyle{primera}{\fancyhf{}\renewcommand{\headrulewidth}{0pt}
  \fancyfoot[L]{\raisebox{-0.32\height}{\includegraphics[height=0.78cm]{logo-gva.png}}\hspace{5mm}%
    \raisebox{-0.32\height}{\includegraphics[height=0.78cm]{logo-vermellar.png}}}
  \fancyfoot[R]{\scriptsize\color{muted}\href{https://creativecommons.org/licenses/by-nc-sa/4.0/deed.ca}{\ccbyncsa}\ \ CC BY-NC-SA 4.0 · fentmates.org\qquad\small\thepage}}
% la resta: el símbol a l'esquerra i el curs a la dreta, amb el títol al centre com abans
\fancyhead[L]{\raisebox{-0.3\height}{\includegraphics[height=0.62cm]{\fm@marca}}}
\fancyhead[R]{\small\bfseries\color{cursfosc}\fm@cursnom}
\renewcommand{\headrule}{\color{curs!45}\hrule height \headrulewidth}
\makeatother

%   \fmcurs{1eso}
% i, al començament del document:
%   \capcalera{Matemàtiques}{1r ESO}{Fitxa d'exercicis · Unitat 2}{Divisibilitat}
% Opcional abans de \input: \def\fmneed{14} (espai mínim abans de cada secció).
% =====================================================================
\makeatletter
\def\fm@logodir{../../logo/pdf}
\providecommand{\fmneed}{10}
\def\fm@cursnom{}

% clau del curs, color i color del text damunt del color
\newcommand{\fm@curs}[3]{\@namedef{fm@color@#1}{#2}\@namedef{fm@text@#1}{#3}}
\fm@curs{1eso}{E6B84A}{navy}
\fm@curs{2eso}{3B82F6}{white}
\fm@curs{3eso}{FF7A50}{white}
\fm@curs{4eso-a}{A78BFA}{white}
\fm@curs{4eso-b}{6D28D9}{white}
\fm@curs{1batx-mi}{9B1C1C}{white}
\fm@curs{1batx-ccss}{EF6B6B}{white}
\fm@curs{2batx-mii}{1E6B47}{white}
\fm@curs{2batx-ccss}{2BB3A3}{white}

\newcommand{\fmcurs}[1]{%
  \@ifundefined{fm@color@#1}{\PackageError{fentmates}{Curs desconegut: #1}{}}{}%
  \edef\fm@tmp{\noexpand\definecolor{curs}{HTML}{\@nameuse{fm@color@#1}}%
               \noexpand\colorlet{curstext}{\@nameuse{fm@text@#1}}}\fm@tmp
  \colorlet{cursfosc}{curs!72!black}%
  \def\fm@logo{\fm@logodir/vertical-#1.pdf}%
  \def\fm@marca{\fm@logodir/marca-#1.pdf}}

% ---------- capçalera de la primera pàgina ----------
% #1 matèria (petit), #2 curs, #3 tipus de document i unitat, #4 títol
\newcommand{\capcalera}[4]{%
  \gdef\fm@cursnom{#2}%
  \par\noindent
  \begin{tcolorbox}[colback=curs,colframe=curs,boxrule=0pt,arc=2.2mm,width=3.5cm,height=2.15cm,
      valign=center,left=3mm,right=2mm,top=1mm,bottom=1mm,nobeforeafter,box align=center]
    \color{curstext}{\footnotesize #1}\par\vspace{1pt}{\fontsize{19}{21}\selectfont\bfseries #2\par}%
  \end{tcolorbox}\hspace{5mm}%
  \begin{minipage}[c]{\dimexpr\linewidth-3.5cm-5mm-2.9cm\relax}
    {\color{cursfosc}\small\bfseries\MakeUppercase{#3}}\par\vspace{2pt}
    {\color{navy}\fontsize{20}{23}\selectfont\bfseries #4\par}
  \end{minipage}\hfill
  \begin{minipage}[c]{2.6cm}\raggedleft\includegraphics[width=2.6cm]{\fm@logo}\end{minipage}\par
  \vspace{2pt}}

% ---------- barres de secció del color del curs ----------
\renewcommand{\seccio}[2]{\par\needspace{\fmneed\baselineskip}\bigskip
  \begin{tcolorbox}[colback=curs,colframe=curs,coltext=curstext,arc=2mm,boxrule=0pt,left=3mm,right=3mm,top=1.2mm,bottom=1.2mm]
  {\large\bfseries #1\quad #2}
  \end{tcolorbox}\vspace{-2pt}}

% ---------- estrelles de dificultat del color del curs (a Batxillerat no se'n posen) ----------
\@ifundefined{estrelles}{}{%
  \renewcommand{\estrelles}[1]{{\color{curs}\ifcase#1\or\faStar\faStar[regular]\faStar[regular]\or\faStar\faStar\faStar[regular]\or\faStar\faStar\faStar\fi}}}

% ---------- capçalera i peu de les pàgines ----------
% primera pàgina: logos institucionals al peu
\fancypagestyle{primera}{\fancyhf{}\renewcommand{\headrulewidth}{0pt}
  \fancyfoot[L]{\raisebox{-0.32\height}{\includegraphics[height=0.78cm]{logo-gva.png}}\hspace{5mm}%
    \raisebox{-0.32\height}{\includegraphics[height=0.78cm]{logo-vermellar.png}}}
  \fancyfoot[R]{\scriptsize\color{muted}\href{https://creativecommons.org/licenses/by-nc-sa/4.0/deed.ca}{\ccbyncsa}\ \ CC BY-NC-SA 4.0 · fentmates.org\qquad\small\thepage}}
% la resta: el símbol a l'esquerra i el curs a la dreta, amb el títol al centre com abans
\fancyhead[L]{\raisebox{-0.3\height}{\includegraphics[height=0.62cm]{\fm@marca}}}
\fancyhead[R]{\small\bfseries\color{cursfosc}\fm@cursnom}
\renewcommand{\headrule}{\color{curs!45}\hrule height \headrulewidth}
\makeatother

%   \fmcurs{1eso}
% i, al començament del document:
%   \capcalera{Matemàtiques}{1r ESO}{Fitxa d'exercicis · Unitat 2}{Divisibilitat}
% Opcional abans de \input: \def\fmneed{14} (espai mínim abans de cada secció).
% =====================================================================
\makeatletter
\def\fm@logodir{../../logo/pdf}
\providecommand{\fmneed}{10}
\def\fm@cursnom{}

% clau del curs, color i color del text damunt del color
\newcommand{\fm@curs}[3]{\@namedef{fm@color@#1}{#2}\@namedef{fm@text@#1}{#3}}
\fm@curs{1eso}{E6B84A}{navy}
\fm@curs{2eso}{3B82F6}{white}
\fm@curs{3eso}{FF7A50}{white}
\fm@curs{4eso-a}{A78BFA}{white}
\fm@curs{4eso-b}{6D28D9}{white}
\fm@curs{1batx-mi}{9B1C1C}{white}
\fm@curs{1batx-ccss}{EF6B6B}{white}
\fm@curs{2batx-mii}{1E6B47}{white}
\fm@curs{2batx-ccss}{2BB3A3}{white}

\newcommand{\fmcurs}[1]{%
  \@ifundefined{fm@color@#1}{\PackageError{fentmates}{Curs desconegut: #1}{}}{}%
  \edef\fm@tmp{\noexpand\definecolor{curs}{HTML}{\@nameuse{fm@color@#1}}%
               \noexpand\colorlet{curstext}{\@nameuse{fm@text@#1}}}\fm@tmp
  \colorlet{cursfosc}{curs!72!black}%
  \def\fm@logo{\fm@logodir/vertical-#1.pdf}%
  \def\fm@marca{\fm@logodir/marca-#1.pdf}}

% ---------- capçalera de la primera pàgina ----------
% #1 matèria (petit), #2 curs, #3 tipus de document i unitat, #4 títol
\newcommand{\capcalera}[4]{%
  \gdef\fm@cursnom{#2}%
  \par\noindent
  \begin{tcolorbox}[colback=curs,colframe=curs,boxrule=0pt,arc=2.2mm,width=3.5cm,height=2.15cm,
      valign=center,left=3mm,right=2mm,top=1mm,bottom=1mm,nobeforeafter,box align=center]
    \color{curstext}{\footnotesize #1}\par\vspace{1pt}{\fontsize{19}{21}\selectfont\bfseries #2\par}%
  \end{tcolorbox}\hspace{5mm}%
  \begin{minipage}[c]{\dimexpr\linewidth-3.5cm-5mm-2.9cm\relax}
    {\color{cursfosc}\small\bfseries\MakeUppercase{#3}}\par\vspace{2pt}
    {\color{navy}\fontsize{20}{23}\selectfont\bfseries #4\par}
  \end{minipage}\hfill
  \begin{minipage}[c]{2.6cm}\raggedleft\includegraphics[width=2.6cm]{\fm@logo}\end{minipage}\par
  \vspace{2pt}}

% ---------- barres de secció del color del curs ----------
\renewcommand{\seccio}[2]{\par\needspace{\fmneed\baselineskip}\bigskip
  \begin{tcolorbox}[colback=curs,colframe=curs,coltext=curstext,arc=2mm,boxrule=0pt,left=3mm,right=3mm,top=1.2mm,bottom=1.2mm]
  {\large\bfseries #1\quad #2}
  \end{tcolorbox}\vspace{-2pt}}

% ---------- estrelles de dificultat del color del curs (a Batxillerat no se'n posen) ----------
\@ifundefined{estrelles}{}{%
  \renewcommand{\estrelles}[1]{{\color{curs}\ifcase#1\or\faStar\faStar[regular]\faStar[regular]\or\faStar\faStar\faStar[regular]\or\faStar\faStar\faStar\fi}}}

% ---------- capçalera i peu de les pàgines ----------
% primera pàgina: logos institucionals al peu
\fancypagestyle{primera}{\fancyhf{}\renewcommand{\headrulewidth}{0pt}
  \fancyfoot[L]{\raisebox{-0.32\height}{\includegraphics[height=0.78cm]{logo-gva.png}}\hspace{5mm}%
    \raisebox{-0.32\height}{\includegraphics[height=0.78cm]{logo-vermellar.png}}}
  \fancyfoot[R]{\scriptsize\color{muted}\href{https://creativecommons.org/licenses/by-nc-sa/4.0/deed.ca}{\ccbyncsa}\ \ CC BY-NC-SA 4.0 · fentmates.org\qquad\small\thepage}}
% la resta: el símbol a l'esquerra i el curs a la dreta, amb el títol al centre com abans
\fancyhead[L]{\raisebox{-0.3\height}{\includegraphics[height=0.62cm]{\fm@marca}}}
\fancyhead[R]{\small\bfseries\color{cursfosc}\fm@cursnom}
\renewcommand{\headrule}{\color{curs!45}\hrule height \headrulewidth}
\makeatother

\fmcurs{2eso}

\begin{document}
\thispagestyle{primera}

\capcalera{Matemáticas}{2\textsuperscript{o} ESO}{Ficha de ejercicios · Unidad 6}{Proporcionalidad y porcentajes}
\vspace{10pt}
Nombre y apellidos: \rule{6.6cm}{0.4pt}\hfill Grupo: \rule{1.6cm}{0.4pt}\hfill Fecha: \rule{2.2cm}{0.4pt}
\vspace{8pt}

\begin{tcolorbox}[colback=soft,colframe=soft,arc=2mm,left=3mm,right=3mm,top=1.5mm,bottom=1.5mm]
\small Ejercicios graduados por niveles de dificultad:\quad \estrelles{1}\ básico \qquad \estrelles{2}\ medio \qquad \estrelles{3}\ avanzado.\\[2pt] Resuélvelos en el cuaderno mostrando todos los pasos. Antes de calcular, decide si la relación es directa o inversa, y comprueba al final que el resultado tiene sentido. Las secciones siguen el orden de los apuntes de la unidad (6.1--6.8).
\end{tcolorbox}

% =====================================================
\seccio{1}{Razones y proporciones}

\ex{1} En un instituto hay 360 alumnos y 24 profesores. Calcula la razón entre alumnos y profesores, simplifícala y explica qué significa.\recorda{La razón entre $a$ y $b$ es el cociente $\frac{a}{b}$. Una proporción es una igualdad de dos razones: $\frac{a}{b}=\frac{c}{d}\iff a\cdot d=b\cdot c$.}

\ex{1} Comprueba si forman proporción:\begin{tasks}(4)
  \task $\frac{3}{4}$ y $\frac{15}{20}$
  \task $\frac{5}{6}$ y $\frac{12}{15}$
  \task $\frac{2{,}5}{2}$ y $\frac{10}{8}$
  \task $\frac{7}{9}$ y $\frac{21}{28}$
\end{tasks}

\ex{2} Calcula el término desconocido:\begin{tasks}(4)
  \task $\frac{x}{12}=\frac{15}{20}$
  \task $\frac{8}{x}=\frac{6}{9}$
  \task $\frac{4{,}5}{6}=\frac{x}{8}$
  \task $\frac{14}{35}=\frac{6}{x}$
\end{tasks}

\ex{2} En un concurso de dibujo se presentan 84 trabajos, de los cuales 35 son de chicos. ¿Cuál es la razón entre trabajos de chicos y de chicas? ¿Y entre los de chicas y el total?

\ex{3} Un batido se hace mezclando leche y fresas en la razón $\frac{5}{2}$. Si ponemos 750 mL de leche, ¿cuántas fresas (en mL) tenemos que poner? ¿Qué cantidad de batido tendremos?

\ex{3} Las edades de dos hermanas están en la razón $\frac{3}{4}$ y entre las dos suman 28 años. ¿Cuántos años tiene cada una?


% =====================================================
\seccio{2}{Magnitudes directamente proporcionales}

\ex{1} Indica qué parejas de magnitudes son directamente proporcionales:\begin{tasks}(2)
  \task El número de kg de manzanas y el precio.
  \task La edad de una persona y su altura.
  \task El lado de un cuadrado y su perímetro.
  \task Los minutos que hablamos por un móvil con cuota fija mensual y el precio de la factura.
\end{tasks}\recorda{Son directamente proporcionales si el doble de una corresponde al doble de la otra, la mitad a la mitad… El cociente entre las cantidades correspondientes es constante.}

\ex{1} Completa la tabla de proporcionalidad directa y calcula su constante:\\[4pt]\begin{tabular}{l|ccccc}Horas & 2 & 5 & 8 &  & 20\\ \hline Km & 7 &  &  & 42 & \end{tabular}

\ex{1} Si 4 entradas de cine cuestan 30 €, ¿cuánto cuestan 7 entradas? Resuélvelo por reducción a la unidad.

\ex{2} Un coche gasta 6,4 L de gasóleo cada 100 km. Resuelve por regla de tres:\begin{tasks}(2)
  \task ¿Cuántos litros gasta en 450 km?
  \task ¿Cuántos km puede hacer con 40 L?
\end{tasks}

\ex{2} Una máquina embotella 1500 botellas en 2 horas y media. ¿Cuántas hará en una jornada de 8 horas? ¿Cuánto tardará en embotellar 2100?

\ex{2} Para hacer una paella para 4 personas hacen falta 320 g de arroz y 1,2 L de caldo. ¿Cuánto hará falta para 10 personas?

\ex{3} En un plano a escala 1:250, una habitación mide 1,6 cm por 1,4 cm. ¿Qué dimensiones reales tiene? ¿Qué superficie tiene, en m²?

\ex{3} Una empresa de alquiler de bicicletas cobra 4 € por hora. Otra cobra 5 € fijos más 3 € por hora. a) Haz una tabla para 1, 2, 3 y 6 horas para cada empresa. b) ¿En qué caso el precio es directamente proporcional al tiempo? c) ¿A partir de cuántas horas es más barata la segunda?


% =====================================================
\seccio{3}{Magnitudes inversamente proporcionales}

\ex{1} Indica si son directamente proporcionales, inversamente proporcionales o ninguna de las dos:\begin{tasks}(2)
  \task El número de pintores y los días que tardan en pintar una casa.
  \task Los litros de agua de un depósito y lo que pesa.
  \task La velocidad de un coche y el tiempo que tarda en un trayecto.
  \task El número de personas de una familia y su gasto en comida.
\end{tasks}\recorda{Son inversamente proporcionales si el doble de una corresponde a la mitad de la otra. El producto de las cantidades correspondientes es constante.}

\ex{1} Completa la tabla de proporcionalidad inversa y calcula su constante:\\[4pt]\begin{tabular}{l|ccccc}Grifos & 2 & 3 & 4 &  & 12\\ \hline Horas & 18 &  &  & 6 & \end{tabular}

\ex{1} Tres máquinas iguales hacen un trabajo en 12 horas. ¿Cuánto tardarían 4 máquinas? Resuélvelo por reducción a la unidad.

\ex{2} Un tren a 120 km/h tarda 3 h 30 min en hacer un trayecto. ¿Cuánto tardaría a 140 km/h?

\ex{2} Un ganadero tiene pienso para alimentar 40 vacas durante 30 días. Si vende 10 vacas, ¿cuántos días le durará el pienso? ¿Y si compra 8 más (en lugar de vender)?

\ex{2} Para llenar un depósito hacen falta 45 cubos de 12 L. ¿Cuántos cubos de 18 L harían falta? ¿De qué capacidad tendrían que ser los cubos para llenarlo con 60 cubos?

\ex{3} Unas obras se tienen que hacer en 18 días y trabajan 10 obreros. Al cabo de 6 días, se van 2 obreros. ¿Cuántos días tardarán los que quedan en acabar la obra?

\ex{3} Decide si cada problema es de proporcionalidad directa o inversa y resuélvelo:\begin{tasks}(1)
  \task Con azulejos de 30 cm de lado necesitamos 24 filas para cubrir una pared. ¿Cuántas filas harían falta con azulejos de 20 cm?
  \task En 5 cajas caben 120 libros. ¿Cuántas cajas iguales hacen falta para 312 libros?
  \task Un grupo de amigos alquila una casa y, si son 6, cada uno paga 85 €. ¿Cuánto pagaría cada uno si fueran 10?
\end{tasks}


% =====================================================
\seccio{4}{Proporcionalidad compuesta}

\ex{1} Para alimentar 6 perros durante 10 días hacen falta 24 kg de pienso. ¿Cuántos kilos hacen falta para 9 perros durante 15 días?\recorda{Compara cada magnitud con la incógnita, suponiendo que las demás no cambian. Si es directa, multiplica por $\frac{\text{nuevo}}{\text{antiguo}}$; si es inversa, por $\frac{\text{antiguo}}{\text{nuevo}}$.}

\ex{1} 4 grifos iguales, abiertos 5 horas, vierten 6000 L. ¿Cuántos litros verterán 6 grifos abiertos 3 horas?

\ex{2} 12 obreros, trabajando 8 horas diarias, hacen una obra en 20 días. ¿Cuántos días tardarán 15 obreros trabajando 8 horas diarias? ¿Y 16 obreros trabajando 6 horas diarias?

\ex{2} Una familia de 4 personas gasta 540 € en comida en 30 días. ¿Cuánto gastará una familia de 6 personas en 20 días?

\ex{2} Para transportar 360 toneladas de tierra, 5 camiones hacen 8 viajes cada uno. ¿Cuántos viajes tendrá que hacer cada camión si son 6 y tienen que transportar 540 toneladas?

\ex{3} En una fábrica, 8 máquinas trabajando 9 horas diarias hacen 3600 piezas en 5 días. ¿Cuántos días necesitarán 6 máquinas trabajando 10 horas diarias para hacer 4800 piezas?

\ex{3} Para arar un campo de 18 ha, 3 tractores trabajan 6 días. ¿Cuántos tractores hacen falta para arar un campo de 30 ha en 5 días?


% =====================================================
\seccio{5}{Repartos proporcionales}

\ex{1} Reparte en partes directamente proporcionales:\begin{tasks}(2)
  \task 840 a 3 y 4
  \task 1200 a 2, 3 y 5
\end{tasks}\recorda{Directo: $k=\dfrac{N}{a+b+c}$ y cada parte es $a\cdot k$, $b\cdot k$… Inverso: se reparte directamente a los inversos $\frac{1}{a},\frac{1}{b}\ldots$ (con el m.c.m., a los numeradores).}

\ex{1} Dos amigos compran un décimo de lotería de 45 €: uno pone 20 € y el otro 25 €. Les toca un premio de 135 €. ¿Cómo se lo tienen que repartir?

\ex{2} Tres socios han invertido 6000 €, 9000 € y 15 000 € en una empresa y han obtenido unos beneficios de 7500 €. ¿Cuánto recibe cada uno?

\ex{2} Reparte en partes inversamente proporcionales:\begin{tasks}(2)
  \task 500 a 2 y 3
  \task 1300 a 2, 3 y 4
\end{tasks}

\ex{2} Una abuela reparte 1170 € entre sus tres nietos de forma directamente proporcional a sus edades: 6, 9 y 11 años. ¿Cuánto recibe cada uno? ¿Y si el reparto fuera inversamente proporcional a 2, 3 y 6?

\ex{3} Un premio de 2800 € se reparte entre tres atletas de forma inversamente proporcional al puesto en que han llegado (1.º, 2.º y 3.º). ¿Cuánto recibe cada uno?

\ex{3} Tres hermanos se reparten una cantidad de forma directamente proporcional a 3, 5 y 7. Al segundo le tocan 60 € más que al primero. ¿Qué cantidad se han repartido? ¿Cuánto recibe cada uno?


% =====================================================
\seccio{6}{Porcentajes: concepto y cálculo}

\ex{1} Completa la tabla:\\[4pt]{\renewcommand{\arraystretch}{2.3}\begin{tabular}[t]{|>{\centering\arraybackslash}m{3.9cm}|>{\centering\arraybackslash}m{3.9cm}|>{\centering\arraybackslash}m{3.9cm}|}\hline \textbf{Porcentaje} & \textbf{Fracción irreducible} & \textbf{Decimal}\\ \hline\hline 30\,\% &  & \\ \hline  & $\dfrac{3}{4}$ & \\ \hline  &  & 0,05\\ \hline 2\,\% &  & \\ \hline  & $\dfrac{1}{8}$ & \\ \hline 150\,\% &  & \\ \hline \end{tabular}}

\ex{1} Calcula mentalmente:\begin{tasks}(4)
  \task 10\,\% de 480
  \task 50\,\% de 76
  \task 25\,\% de 360
  \task 5\,\% de 240
  \task 20\,\% de 45
  \task 1\,\% de 3500
  \task 15\,\% de 60
  \task 75\,\% de 80
\end{tasks}\recorda{El 10 \% es dividir entre 10; el 5 \%, la mitad del 10 \%; el 50 \%, la mitad; el 25 \%, la cuarta parte; el 20 \%, la quinta parte.}

\ex{1} Calcula:\begin{tasks}(3)
  \task 12\,\% de 350
  \task 35\,\% de 84
  \task 4\,\% de 1250
  \task 7{,}5\,\% de 240
  \task 0{,}5\,\% de 6000
  \task 120\,\% de 65
\end{tasks}

\ex{2} En un instituto de 850 alumnos, el 46 \% son chicos, el 18 \% va a la escuela de música y el 6 \% es de otros países. ¿Cuántos alumnos hay en cada caso?

\ex{2} El cuerpo humano es aproximadamente un 65 \% agua. ¿Cuántos kilos de agua tiene una persona de 58 kg? ¿Y una de 72 kg?

\ex{3} De un depósito de 2000 L, el lunes se gasta el 30 \% y el martes, el 25 \% de lo que quedaba. ¿Cuántos litros quedan? ¿Qué porcentaje del depósito se ha gastado?


% =====================================================
\seccio{7}{¿Qué porcentaje representa? Cálculo del total}

\ex{1} ¿Qué porcentaje representa…\begin{tasks}(2)
  \task 12 de 50?
  \task 27 de 36?
  \task 45 de 300?
  \task 90 de 72?
\end{tasks}\recorda{$p\,\%=\frac{\text{parte}}{\text{total}}\cdot100$. \quad Total $=$ parte $:\ \frac{p}{100}$.}

\ex{1} Calcula la cantidad total si:\begin{tasks}(2)
  \task el 20 \% es 36
  \task el 15 \% es 63
  \task el 75 \% es 240
  \task el 4 \% es 18
\end{tasks}

\ex{2} En una clase de 28 alumnos han aprobado 21. ¿Qué porcentaje ha aprobado? ¿Y cuál ha suspendido?

\ex{2} En una votación, una propuesta ha obtenido 182 votos, que son el 35 \% de los votos emitidos. ¿Cuántas personas han votado?

\ex{2} Una jugadora de baloncesto ha hecho 15 tiros de 3 puntos y ha acertado 6. Un compañero ha hecho 20 y ha acertado 7. ¿Quién tiene mejor porcentaje de acierto?

\ex{3} De una pizza nos hemos comido el 62,5 \% y quedan 3 porciones iguales. ¿En cuántas porciones estaba cortada? ¿Qué porcentaje representa cada porción?

\ex{3} Una tienda ha vendido este mes 252 bicicletas, que son el 120 \% de las que vendió el mes pasado. ¿Cuántas vendió el mes pasado? ¿Cuántas más ha vendido este mes?


% =====================================================
\seccio{8}{Aumentos y disminuciones porcentuales}

\ex{1} Escribe el índice de variación de cada caso:\begin{tasks}(3)
  \task aumento del 25 \%
  \task disminución del 10 \%
  \task aumento del 4 \%
  \task rebaja del 35 \%
  \task IVA del 21 \%
  \task aumento del 100 \%
\end{tasks}\recorda{Aumento del $p\,\%$: índice $1+\frac{p}{100}$. Disminución del $p\,\%$: índice $1-\frac{p}{100}$. \ Final $=$ inicial $\cdot$ índice; \ inicial $=$ final $:$ índice.}

\ex{1} Calcula el precio final:\begin{tasks}(2)
  \task Una chaqueta de 90 € rebajada un 30 \%.
  \task Un libro de 22 € con un 4 \% de IVA.
  \task Un sueldo de 1450 € que sube un 3 \%.
  \task Una tableta de 340 € con un 15 \% de descuento.
\end{tasks}

\ex{2} Un coche valía 18 500 € y ha bajado un 12 \%. ¿Cuánto vale ahora? ¿Cuántos euros ha bajado?

\ex{2} Calcula la cantidad inicial:\begin{tasks}(1)
  \task Después de una rebaja del 20 \%, unas botas cuestan 68 €.
  \task Con el IVA del 21 \%, una factura es de 157,30 €.
  \task Después de aumentar un 6 \%, un pueblo tiene 3180 habitantes.
\end{tasks}

\ex{2} Calcula el porcentaje de variación:\begin{tasks}(1)
  \task Una camiseta pasa de 25 € a 20 €.
  \task Una ciudad pasa de 40 000 a 43 000 habitantes.
  \task El precio del aceite pasa de 6,40 € a 8 € el litro.
\end{tasks}

\ex{3} El precio de una consola de 400 € sube un 10 \% y, después, baja un 10 \%. ¿Cuánto vale al final? ¿Cuál es la variación porcentual total?

\ex{3} En unas rebajas hacen un 40 \% de descuento y, a la hora de pagar, un 10 \% más sobre el precio rebajado. a) ¿Es lo mismo que un 50 \% de descuento? b) Si hemos pagado 37,80 €, ¿cuánto costaba antes?

\ex{3} Un comerciante compra una bicicleta por 250 €, la pone a la venta ganando un 40 \% y, como no la vende, después la rebaja un 20 \%. a) ¿A qué precio la vende finalmente? b) ¿Qué porcentaje gana respecto del precio de compra?

\vfill
\begin{center}\color{muted}\small\faCheckCircle[regular]\ \ ¡Ánimo! Comprueba tus resultados con las herramientas interactivas de los apuntes de la unidad.\end{center}

\end{document}
